\documentclass[UTF8,a4paper,11pt]{ctexart}
\usepackage[margin=25mm]{geometry}
\usepackage{graphicx,booktabs,tabularx,array,float,enumitem}
\usepackage[hidelinks]{hyperref}
\usepackage{xcolor}
\graphicspath{{../figures/}}
\setlength{\parindent}{2em}
\setlength{\parskip}{0.35em}
\linespread{1.35}
\newcolumntype{Y}{>{\raggedright\arraybackslash}X}
\hypersetup{pdftitle={Shelfwise库存与缺货预警系统——中文逻辑对照稿},pdfauthor={作者姓名（占位）}}
\newcommand{\fig}[3]{\begin{figure}[H]\centering\includegraphics[width=\textwidth,height=0.42\textheight,keepaspectratio]{#1}\caption{#2}\label{#3}\end{figure}}
\begin{document}
\begin{titlepage}
\centering
{\zihao{-4}学校名称（占位）\\院系名称（占位）\par}
\vfill
{\zihao{2}\bfseries 面向小型超市的库存与缺货预警系统\par}
\vspace{1em}
{\zihao{3}英文论文逻辑对照稿\par}
\vspace{3em}
{\zihao{4}作者姓名（占位）\\学号（占位）\\指导教师（占位）\par}
\vfill
{\zihao{-4}本稿用于中文阅读和论证检查，不替代英文提交版本。\\2026年10月\par}
\end{titlepage}
\pagenumbering{Roman}
\section*{说明}
本稿与英文论文使用同一套系统、代码、测试和实验数据，目的在于帮助读者用中文检查论文的逻辑链条。中文稿不重新提出一套不同的方法，也不增加英文论文中没有验证的结论。作者、学校、院系、学位、导师和学号等信息暂时保留占位符。

论文的主线可以概括为：小型超市需要的不只是一个“当前库存数”，还需要知道库存为什么变化、谁进行了操作，以及哪些商品需要处理。因此，系统把库存数量、库存流水、预警生命周期和人工确认分开建模，再通过数据库事务、行锁、幂等键和可失效缓存保证这几类信息不会互相矛盾。论文后半部分不把系统功能直接当成研究结论，而是使用测试和实验逐项验证前面的设计目标。

\setcounter{tocdepth}{1}
\tableofcontents
\clearpage
\pagenumbering{arabic}
\section{问题、目标与研究问题}
\subsection{问题背景}
小型超市的库存变化频繁而且来源多样：早上的送货会增加库存，销售或出库会减少库存，人工盘点可能发现差异，损耗又可能需要单独调整。如果系统只保存一个可以直接编辑的数量，工作人员无法在事后区分“收到了一批货”“卖出了一件商品”“盘点修正了数量”或者“商品损坏”。数量即使暂时正确，解释也已经丢失。

缺货预警还有一个容易被忽略的逻辑问题。店长点击“已查看”并不会让商品重新上架；相反，补货成功也不一定意味着有人点击了确认按钮。如果把确认动作直接当成预警解决，系统会把人的处理行为和库存事实混在一起。因此，本项目把“库存是否恢复”和“谁查看过预警”建模为两个相互独立的维度。

\subsection{研究目标}
项目目标是实现并验证一个单店库存与缺货预警系统 Shelfwise。系统需要满足以下条件：库存变化必须能够追溯；并发或重复请求不能把库存扣成负数或重复记账；预警状态应该能解释库存状态和人工处理状态；Redis 只能作为查询加速手段，不能成为库存正确性的前提；其他开发者能够从干净检出开始复现主要检查和论文构建。

\subsection{研究问题}
\begin{table}[H]\centering\small
\caption{研究问题与证据对应关系}
\begin{tabularx}{\textwidth}{p{0.12\textwidth}YY}\toprule
问题 & 关注点 & 主要证据\\\midrule
RQ1 & 重复提交和并发出库下，库存是否保持非负并保留可归因历史？ & MySQL/Redis 集成测试、库存流水、并发场景\\
RQ2 & 缺货、补货和人工确认能否表示为可观察的不同事件？ & 预警状态迁移、阈值测试、预警界面和确认测试\\
RQ3 & Redis 缓存对本地预警查询有什么影响，Redis 故障时结果是否保持？ & 交替缓存实验、响应等价性、Redis 停机恢复实验\\
RQ4 & 其他开发者能否复现系统检查和论文产物？ & 版本矩阵、迁移脚本、CI、构建和证据脚本\\\bottomrule
\end{tabularx}
\end{table}

这些问题有意保持为软件工程问题。RQ1不是证明所有硬件故障下的绝对正确性；RQ2不是用户满意度调查；RQ3不是生产环境 SLA；RQ4也不是保证任何机器都能零配置运行。每个问题都限定了研究对象、实验环境和证据类型，从而避免让论文结论超过实现和实验实际支持的范围。

\subsection{范围边界}
系统面向一个小型超市，支持管理员、店长和库存员三个角色。系统包含商品、分类、供应商、库存事务、库存流水、缺货预警、预警确认、报表和基础设置。系统不包含多店租户、在线支付、财务核算、采购结算、自动采购订单和需求预测；库存使用整数单位，不处理小数商品和批次有效期。

库存价值仅按当前数量乘以当前商品单价计算，是一个展示估计，不是会计成本或历史估值。补货阈值由店长维护，系统验证阈值之间的数值关系，但不声称能够自动推导最优阈值。

\section{需求与验收标准}
\subsection{角色职责}
管理员负责产品、分类、供应商、用户和设置；店长查看库存和报表、调整安全库存与补货阈值、查看并确认预警；库存员登记入库、出库、调整和盘点。三类用户都可以查询库存和库存流水，但服务器仍然逐请求检查权限，前端隐藏菜单不是唯一的安全边界。

\begin{table}[H]\centering\small
\caption{角色权限矩阵}
\begin{tabularx}{\textwidth}{Yccc}\toprule
操作 & 管理员 & 店长 & 库存员\\\midrule
查询库存和库存流水 & 允许 & 允许 & 允许\\
登记入库、出库、调整、盘点 & 允许 & 不允许 & 允许\\
维护商品、分类、供应商 & 允许 & 不允许 & 不允许\\
调整安全库存和补货阈值 & 允许 & 允许 & 不允许\\
查看和确认预警 & 允许 & 允许 & 不允许\\
查看报表 & 允许 & 允许 & 不允许\\
管理用户和系统设置 & 允许 & 不允许 & 不允许\\\bottomrule
\end{tabularx}
\end{table}

\subsection{库存不变量}
库存事务使用如下关系：
\[
q_{new}=q_{old}+\Delta,
\qquad 0\leq q_{new}\leq 1{,}000{,}000{,}000。
\]
入库和出库的输入数量为正数，由服务层决定符号；调整使用带符号数量；盘点使用盘点后的绝对数量，系统在持有产品行锁时计算差值。库存数量不能通过商品编辑接口直接覆盖。

每条成功流水至少包含产品、操作人、事务类型、变更前数量、变更量、变更后数量、原因、幂等键和服务器时间。原因必须存在，但原因文字本身不能证明现实中的送货或盘点一定发生，因此论文只把它称为可追溯记录，不称为不可篡改的法律证据。

\subsection{预警语义}
数量为 0 时生成 OUT；数量大于 0 且小于等于补货阈值时生成 LOW；数量高于补货阈值时处于健康状态。缺货恢复后关闭原缺货事件，并创建 24 小时有效的 RESTOCKED 事件。确认预警只增加确认人和确认时间，不改变数量，也不自动关闭预警。状态变化、阈值修改和归档操作都需要重新计算预警。

\section{系统设计}
\subsection{整体架构}
系统是单体模块化架构：Vue 3 浏览器端负责界面、路由、表单和状态展示；Spring Boot 负责 REST API、认证、权限、业务事务和校验；MySQL 保存持久化数据；Redis 缓存预警查询结果。MySQL 是最终事实来源，Redis 故障时查询可以回退到 MySQL。

\fig{architecture.png}{Shelfwise 系统运行架构。MySQL 保存持久状态，Redis 只负责可选查询缓存。}{fig:cn-architecture}

选择模块化单体而不是拆分多个服务，原因是项目范围只有一个门店，库存和预警又共享同一个事务边界。拆分服务会引入消息延迟、重复消费和跨服务一致性问题，但当前没有证据表明系统需要这种复杂度。

\subsection{数据库模型}
核心实体包括用户、商品、分类、供应商、库存事务、预警、设置和缓存版本。商品表保存当前数量，库存事务表保存变化解释，二者在同一个事务中更新。事务表使用“操作人 + 幂等键”唯一约束，商品 SKU 和条码使用唯一约束。商品有历史记录后只能归档，不能直接删除；归档不会伪造出库流水。

\fig{erd.png}{核心实体关系。库存流水和预警均引用商品，操作人负责提供归因和确认信息。}{fig:cn-erd}

\subsection{库存事务算法}
一次入库、出库或盘点请求依次经过：身份和角色校验、输入校验、产品行锁、幂等键检查、库存边界检查、库存与流水写入、预警更新以及缓存版本递增。任何一步失败，整个数据库事务回滚。

行锁解决并发出库问题。例如库存只有 1 件时，两个并发出库请求不能都读取到旧值 1 并同时成功。第一个请求提交后，第二个请求读取到最新数量 0 并被拒绝。幂等键解决网络重试问题：相同操作重试时返回原流水；相同键但内容不同则返回冲突。

\fig{stock-workflow.png}{库存事务和预警更新流程。库存、流水、预警版本在同一事务边界内处理。}{fig:cn-workflow}

\subsection{预警生命周期与缓存}
预警生命周期和人工确认分开存储。反复查询同一缺货事件不会产生无限重复预警；严重程度变化会关闭旧事件并开启新事件；库存恢复会产生短期 RESTOCKED 事件。Redis 键包含数据库缓存版本、预警类型、状态、页码和页大小。库存、阈值或确认发生提交时，数据库版本递增，后续读取自然进入新命名空间。旧缓存由 TTL 回收，而不是清空整个 Redis。

\section{系统实现}
\subsection{认证和权限}
系统使用服务端 Session、CSRF Token 和 BCrypt 密码哈希。浏览器先获取 CSRF Token，再提交登录或其他写操作。用户状态和角色在请求期间从 MySQL 重新确认，因此暂停账号或修改角色后不依赖用户主动退出登录才能生效。演示账号和密码来自被 Git 忽略的本地环境文件。

\subsection{库存和预警界面}
库存页面支持按名称、SKU、条码、分类和库存状态搜索，并提供分页和稳定排序。库存事务表单支持收货、出库、调整和盘点，并要求填写原因。预警页面同时展示商品、当前数量、阈值、首次发现时间和确认状态。报表页面把当前库存数量、当前价格和交易日期聚合成可读指标，但明确说明这些数据不是会计估值或销售预测。

\fig{ui-dashboard.png}{系统首页仪表盘，展示当前目录、库存预警、库存估值估计和记录的事务数量。}{fig:cn-dashboard}

\fig{ui-stock-form.png}{库存事务表单。表单提交的是操作意图，数量前后值由服务器在事务中计算。}{fig:cn-stock}

\subsection{实现中的修正}
编译通过并不等于系统行为正确。开发过程中发现并修正了 Vite 依赖缓存、ECharts 渲染器注册、库存流水懒加载、Java 时间序列化和环境变量映射等问题。这些问题说明系统必须同时经过源代码检查、数据库接口调用和浏览器运行检查。

\section{测试与实验}
\subsection{测试结果}
后端使用真实 MySQL 和 Redis 容器运行 27 项集成测试，覆盖权限、CSRF、阈值边界、重复提交、并发出库、归档、预警确认、管理员保护和输入校验。前端有 3 项库存状态单元测试。浏览器脚本完成 14 项检查，包括登录、库存搜索、库存登记、预警确认、角色跳转和移动端布局，未发现 JavaScript 运行时错误。

\begin{table}[H]\centering\small
\caption{主要验证结果}
\begin{tabularx}{\textwidth}{p{0.28\textwidth}p{0.17\textwidth}Y}\toprule
验证层 & 结果 & 说明\\\midrule
后端集成测试 & 27 项通过 & 使用真实 MySQL/Redis，覆盖事务、权限、并发和预警\\
前端单元测试 & 3 项通过 & 覆盖零库存、等于阈值和健康库存\\
浏览器流程 & 14 项通过 & 页面渲染、库存登记、预警确认和移动端检查\\
Redis 故障恢复 & 通过 & Redis 停止期间仍返回同一数据库结果\\
应用重启持久性 & 通过 & 重启前后库存与流水指纹一致\\\bottomrule
\end{tabularx}
\end{table}

\subsection{缓存实验}
缓存实验使用同一个认证用户、同一个接口、同一个页大小和同一个演示数据集。缓存路径和直接查询路径各采集 600 个请求样本。缓存路径的中位延迟约为 5.33 ms，直接查询约为 6.56 ms；缓存和非缓存响应经过 JSON 等价性比较。该结果只能说明在本地小数据集上缓存减少了一部分查询工作，不能推导出生产环境 SLA，也不能证明所有小型超市都必须使用 Redis。

\subsection{故障与持久性}
停止 Redis 后，预警接口仍通过 MySQL 返回相同的商品、数量和确认信息；Redis 恢复后查询结果保持一致。应用重启后，商品数量和库存流水的有序指纹不变。这两个实验支持“Redis 不是写入正确性的前提”和“正常应用重启不丢失库存历史”这两个有限结论，但不等于已经完成磁盘损坏恢复或生产灾备验证。

\section{讨论、局限与结论}
项目证明了一个单店库存演示系统可以把库存余额、流水历史、预警生命周期和人工确认组织在同一套可复现实现中。行锁、幂等键和数据库约束共同保护库存不变量；预警确认不再冒充补货；缓存版本让查询加速不改变 MySQL 的权威地位。

项目没有证明降低了真实缺货损失，也没有做真实门店用户研究。库存使用整数单位，没有批次和有效期，阈值由人工配置，没有多店部署、自动采购、财务核算和需求预测。下一步应优先进行门店任务观察、安全和备份恢复评估、分数单位和批次支持、POS 接入以及更大规模的并发实验，而不是继续堆叠新的技术组件。

因此，论文的结论不是“系统已经适用于所有超市”，而是：在明确的单店范围和本地验证环境内，Shelfwise 实现了可追溯库存事务、可解释预警和可回退缓存路径，并且这些结论可以由源代码、数据库测试、浏览器证据和实验数据逐项复核。

\section*{阅读对照建议}
阅读英文论文时，可以依次检查四条线：第一，绪论中的研究问题是否在需求和测试中有对应物；第二，设计章节中的库存不变量是否在实现和集成测试中落地；第三，Redis 性能结果是否同时报告了缓存一致性和故障回退；第四，结论是否只复述已经验证的内容，并主动列出真实部署仍然缺少的证据。若这四条线闭合，论文的整体论证就成立。

\end{document}
